\begin{center}[9.]\par\medskip\textit{Induction by Motion.}\end{center}
I.\quad Take any point in space with coordinates $x,y,z$. Let $V$ denote the solid angle that a current $S'$ subtends at that point.

$V=\text{Const.}$ thus determines a surface whose tangent plane and normal are what \textsc{Ampère} calls the Plan directeur. and Directrice.

II.\quad In every moving material molecule $\mu$ with velocity components $\frac{\mathrm dx}{\mathrm dt},\frac{\mathrm dy}{\mathrm dt},\frac{\mathrm dz}{\mathrm dt}$, an electrodynamic \textit{force} arises. Let $\xi,\eta,\zeta$ be its components. Then
\[
\begin{aligned}
\xi&=\left(\frac{\mathrm dV}{\mathrm dy}\,.\,\frac{\mathrm dz}{\mathrm dt}-\frac{\mathrm dV}{\mathrm dz}\,.\,\frac{\mathrm dy}{\mathrm dt}\right)\mu\\[4pt]
\eta&=\left(\frac{\mathrm dV}{\mathrm dz}\,.\,\frac{\mathrm dx}{\mathrm dt}-\frac{\mathrm dV}{\mathrm dx}\,.\,\frac{\mathrm dz}{\mathrm dt}\right)\mu\\[4pt]
\zeta&=\left(\frac{\mathrm dV}{\mathrm dx}\,.\,\frac{\mathrm dy}{\mathrm dt}-\frac{\mathrm dV}{\mathrm dy}\,.\,\frac{\mathrm dx}{\mathrm dt}\right)\mu
\end{aligned}
\]

\SourceLeaf{618}III.\quad If, on the other hand, an element $\mathrm ds$ of a current-carrying body passes through that point, $S$ acts on the body. If the component forces are $\xi,\eta,\zeta$ and $i$ is the current intensity, then
\[
\begin{aligned}
\xi&=\left(\frac{\mathrm dV}{\mathrm dy}\,.\,\frac{\mathrm dz}{\mathrm ds}-\frac{\mathrm dV}{\mathrm dz}\,.\,\frac{\mathrm dy}{\mathrm ds}\right)i\,\mathrm ds\\
&\text{etc.}
\end{aligned}
\]

Also,
\[
\begin{aligned}
V&=\int\frac{z'(x'\mathrm dy'-y'\mathrm dx')}{r'(x'x'+y'y')}=\int\frac{x'(y'\mathrm dz'-z'\mathrm dy')}{r'(y'y'+z'z')}\\
&=\int\frac{y'(z'\mathrm dx'-x'\mathrm dz')}{r'(z'z'+x'x')}
\end{aligned}
\]
where $x',y',z'$ refer to the \textit{acting} current $S'$. It follows that
\[
\frac{\mathrm dV}{\mathrm dz}=\int\frac{x'\mathrm dy'-y'\mathrm dx'}{r'^3}
\]
which is \textsc{Ampère}’s expression. For the indefinite integral one has
\[
\int\left\{\frac{z(x\mathrm dy-y\mathrm dx)}{r(xx+yy)}-\frac{y(z\mathrm dx-x\mathrm dz)}{r(xx+zz)}\right\}=\text{Arc.}\,\operatorname{tg}\frac{yz}{xr}
\]
This is proved most simply by writing it in the form
\[
\frac{-yzrr\,\mathrm dx-xyzx\,\mathrm dx+xzrr\,\mathrm dy-xyzy\,\mathrm dy+xyrr\,\mathrm dz-xyzz\,\mathrm dz}{r(xx+yy)(xx+zz)}
\]

\vspace{4mm}
The total current-producing force arising from the \textit{formation of a current} at any point in space, for example at the point whose coordinates are all $=0$, is resolved into three component forces, namely
\[
\begin{aligned}
X&=\int\frac{\mathrm dx'}{r}&\quad\text{so }\frac{\mathrm dV}{\mathrm dx}&=\frac{\mathrm dZ}{\mathrm dy}-\frac{\mathrm dY}{\mathrm dz}\\[4pt]
Y&=\int\frac{\mathrm dy'}{r}&\frac{\mathrm dV}{\mathrm dy}&=\frac{\mathrm dX}{\mathrm dz}-\frac{\mathrm dZ}{\mathrm dx}\\[4pt]
Z&=\int\frac{\mathrm dz'}{r}&\frac{\mathrm dV}{\mathrm dz}&=\frac{\mathrm dY}{\mathrm dx}-\frac{\mathrm dX}{\mathrm dy}
\end{aligned}
\]

(February 4.) [Later addition:] (1836 April 7.)

(For a plane curve, $z'$ is constant, so $Z=0$. Let
\[
\int r\frac{(x'-x)\mathrm dy'-(y'-y)\mathrm dx'}{(x'-x)^2+(y'-y)^2}=U
\]
\[
U=\int r\,\mathrm d'\lambda=(z'-z)\int\frac{\mathrm d'\lambda}{\sin\beta}\quad\text{if}\quad z'-z=r\sin\beta,\quad\frac{y'-y}{x'-x}=\operatorname{tg}\lambda
\]
\SourceLeaf{619}then one obtains
\[
\begin{aligned}
X&=\frac{\mathrm dU}{\mathrm dy}\\[4pt]
Y&=-\frac{\mathrm dU}{\mathrm dx}
\end{aligned}
\]

Thus the direction of the force with components $X,Y,Z$ always lies in the surface on which $U=\text{const.}$ and at the same time in the plane on which $z$ is constant. This line therefore closes on itself.)

One has
\[
\begin{aligned}
\xi&=\frac{\mathrm dV}{\mathrm dy}\,.\,\mathrm dz-\frac{\mathrm dV}{\mathrm dz}\,.\,\mathrm dy\\[4pt]
&=\int'\frac{\begin{gathered}((z'-z)\mathrm dx'-(x'-x)\mathrm dz')\mathrm dz\\-((x'-x)\mathrm dy'-(y'-y)\mathrm dx')\mathrm dy\end{gathered}}{r^3}\\[4pt]
&=\int'\left\{\begin{gathered}\frac{\mathrm dx'[(x'-x)\mathrm dx+(y'-y)\mathrm dy+(z'-z)\mathrm dz]}{r^3}\\-\frac{(x'-x)(\mathrm dx\,\mathrm dx'+\mathrm dy\,\mathrm dy'+\mathrm dz\,\mathrm dz')}{r^3}\end{gathered}\right\}
\end{aligned}
\]
hence
\[
\begin{aligned}
\xi\delta x+\eta\delta y+\zeta\delta z&=-\int'\frac{\mathrm dr}{rr}(\mathrm dx'\,.\,\delta x+\mathrm dy'\,.\,\delta y+\mathrm dz'\,.\,\delta z)\\
&\qquad+\int'\frac{\delta r}{rr}(\mathrm dx'\,.\,\mathrm dx+\mathrm dy'\,.\,\mathrm dy+\mathrm dz'\,.\,\mathrm dz)\\
&=-\int'\delta\,.\,\frac{\mathrm dx\,.\,\mathrm dx'+\mathrm dy\,.\,\mathrm dy'+\mathrm dz\,.\,\mathrm dz'}{r}\\
&\qquad+\int'\left(\mathrm dx'\,.\,\frac{\delta\mathrm dx}{r}+\mathrm dy'\,.\,\frac{\delta\mathrm dy}{r}+\mathrm dz'\,.\,\frac{\delta\mathrm dz}{r}\right)\\
&\qquad-\int'\left(\mathrm dx'\,.\,\frac{\delta x\,.\,\mathrm dr}{rr}+\mathrm dy'\,.\,\frac{\delta y\,.\,\mathrm dr}{rr}+\mathrm dz'\,.\,\frac{\delta z\,.\,\mathrm dr}{rr}\right)
\end{aligned}
\]
The last two parts become
\[
\int'\mathrm dx'\,.\,\mathrm d\frac{\delta x}{r}+\mathrm dy'\,.\,\mathrm d\frac{\delta y}{r}+\mathrm dz'\,.\,\mathrm d\frac{\delta z}{r}
\]
which vanishes when integrated over all of $S$. Thus
\[
\int\xi\delta x+\eta\delta y+\zeta\,\mathrm dz=-\delta\int\int'\frac{\mathrm dx\,.\,\mathrm dx'+\mathrm dy\,.\,\mathrm dy'+\mathrm dz\,.\,\mathrm dz'}{r}
\]

A shorter proof proceeds as follows. One has
\[
\begin{aligned}
\text{1)}\qquad\delta\int\int\rho\,\mathrm d\mathrm d'\rho&=\int\int\delta\rho\,.\,\mathrm d\mathrm d'\rho+\int\int\rho\delta\mathrm d\mathrm d'\rho\\[4pt]
\text{2)}\qquad\int\int\rho\,\mathrm d'\mathrm d\delta\rho+\int\int\mathrm d'\rho\,.\,\mathrm d\delta\rho&=\int\int\mathrm d'(\rho\,\mathrm d\delta\rho)=0\\[4pt]
\text{3)}\qquad0=\int\int\mathrm d(\delta\rho\,.\,\mathrm d'\rho)&=\int\int\mathrm d'\rho\,.\,\mathrm d\delta\rho+\int\int\mathrm d\mathrm d'\rho\,.\,\delta\rho
\end{aligned}
\]
\SourceLeaf{620}so by addition
\[
\delta\int\int\rho\,\mathrm d\mathrm d'\rho=2\int\int\delta\rho\,.\,\mathrm d\mathrm d'\rho=\int\int\delta r\,.\,\frac{\mathrm d\mathrm d'\rho}{\rho}
\]

The double integral under the variation sign can be expressed in various ways
\[
\begin{aligned}
\int\int\rho\,\mathrm d\mathrm d'\rho&=-\int\int\mathrm d\rho\,.\,\mathrm d'\rho\\
&=\int\int\frac{1}{4r}\cos\theta\,.\,\cos\theta'\,.\,\mathrm ds\,.\,\mathrm ds'\\
&=\int\int\frac{1}{4r}(\cos\theta\,.\,\cos\theta'+\sin\theta\sin\theta'\cos\omega)\mathrm ds\,.\,\mathrm ds'\\
&=\frac{1}{m}\,.\,\frac{1}{2-m}\int\int\rho^{2-m}\mathrm d\mathrm d'\rho^m\\
&=-\frac{1}{mm}\int\int\rho^{2-2m}\mathrm d\rho^m\,.\,\mathrm d'\rho^m
\end{aligned}
\]

